\subsection{Congruence descent and the precise small-level obstruction}

\begin{result}{D.31}{Proved}{The full GT action descends through the actual congruence quotients}
For every stable hyperbolic type $(g,n)$, put
\[
 \Pi_{g,n}=\widehat{\pi_1(S_{g,n})},\qquad
 K_{g,n}=\ker\bigl(\widehat\Gamma_{g,n}
                 \longrightarrow\operatorname{Out}(\Pi_{g,n})\bigr).
                                                               \tag{D.40}
\]
Every parameter of $\widehat{\mathrm{GT}}$ in E.39 preserves
$K_{g,n}$.  Its full exit-system action descends to the
procongruence exit categories defined by the inverse system
of the actual finite ambient transporter quotients (E.82),
with the index restricted to congruence normal subgroups.
The descended action commutes with relabeling, nodal clutching
and forgetting, including their stated coherence maps.  Its
genus-zero restriction is the original GT parameter.
The assertion includes closed surfaces of genus at least two;
the unstable unmarked torus is outside this hyperbolic range.
\end{result}
\begin{proof}
We give the full Birman comparison used to recognize the
kernel in (D.40).  For a stable hyperbolic surface there is
an exact sequence of full completions
\[
 1\longrightarrow\Pi_{g,n}\longrightarrow
 \widehat\Gamma_{g,n+1}\longrightarrow
 \widehat\Gamma_{g,n}\longrightarrow1.                       \tag{D.41}
\]
For $n>0$ the proof is the characteristic-finite-quotient
conjugation argument of E.33.  We spell out the closed case
as well.  The profinite group of a closed surface of genus
$g\geq2$ is center-free.  Indeed suppose that a central
$z\ne1$ is outside an open normal subgroup $\widehat H$,
and choose a discrete surface-group element $a$ representing
$z\widehat H$.  Conjugation by $a$ on $\widehat H$ is inner,
because it differs from conjugation by $z$ by an element of
$\widehat H$.  It is therefore the identity on
$H_1(H;\widehat{\mathbb Z})$ and on the integral homology of
the corresponding finite closed-surface cover $S_H$.
But $aH$ is a nonidentity deck transformation.  Lift a finite
cell decomposition from the base surface: the deck
transformation freely permutes the cells above every base
cell, so its alternating cellular trace is zero.  Its
homological trace is instead $2-2g(S_H)$, since it preserves
orientation and acts trivially on $H_1$.  This is nonzero,
as $g(S_H)\geq2$, a contradiction.

For either the free or the closed surface group, characteristic
finite-index subgroups are cofinal.  Conjugation in the
discrete Birman extension therefore extends to an action of
the full completed total group on $\Pi_{g,n}$.  On the
completed point-pushing group it is its inner action.  The
center-free result makes the map of that completed kernel
into the total completion injective.  Right exactness of
completion proves (D.41).  The induced outer conjugation
map of its quotient is exactly the monodromy in (D.40),
since this equality holds in the discrete evaluation
fibration and both actions are continuous.

Let $F$ be one of the full-system equivalences of E.39 and
write $\alpha$ for its automorphism of the base group in
(D.41).  Its compatibility with the actual forgetting map
gives an automorphism of that extension.  If the selected
basepoint comparisons conjugate the forgetting square,
lift that conjugator through the surjection (D.41) and
adjust the total-group coordinates; this makes that square
commute literally.  Its restriction to the kernel is a
continuous automorphism $\beta$ of $\Pi_{g,n}$.  Conjugation
in the extension gives the actual identity
\[
 \rho(\alpha(x))=[\beta]\rho(x)[\beta]^{-1}
          \quad\text{in }\operatorname{Out}(\Pi_{g,n}),       \tag{D.42}
\]
where $\rho$ is (D.40).  Applying the same equation to
$F^{-1}$ proves $\alpha(K_{g,n})=K_{g,n}$.  This statement
is unchanged by any of the indicated inner-coordinate
choices.  No congruence-subgroup-property assumption was
used.

We next prove descent on the arrows, including their
possible strict-completion ambiguity.  Fix the actual
full-smoothing crossing $j_\sigma=[1]:\sigma\to\varnothing$
at every object.  By E.29 there is $u_\sigma\in
\widehat\Gamma$ with $F(j_\sigma)=u_\sigma j_\sigma$.
Its ambiguity is exactly right multiplication by
$\widehat\Delta_\sigma$.  The left stabilizer of this
crossing gives
$\alpha(\widehat\Delta_\sigma)=
u_\sigma\widehat\Delta_\sigma u_\sigma^{-1}$.
For an arrow represented by $g:\sigma\to\tau$,
functoriality of $j_\tau g=g j_\sigma$ gives, in every
finite ambient transporter category, the coordinate
\[
       [g]\longmapsto
       [u_\tau^{-1}\alpha(g)u_\sigma].                       \tag{D.43}
\]
To see this without assuming injectivity of the strict
completed Hom into its open composite, project to (E.82).
The image of $F(j_\tau)$ is represented by $q_N(u_\tau)$;
cancel that invertible finite-group element from the
equality of composites.  The remaining equality is exactly
(D.43).  The displayed stabilizer identity proves
independence from a source twist or from the allowed
crossing representative.

For a congruence normal subgroup $N$, its inverse image
$\alpha^{-1}(N)$ is again congruence by (D.42).  Thus
(D.43) factors through a finite ambient congruence
quotient of the source.  The actual full functor ensures
that its value belongs to the indicated transporter
subset.  Its formulas commute with refinement of $N$,
and $F^{-1}$ supplies their inverse.  Their limits give
the asserted continuous equivalences of procongruence
exit categories.  Multiplication of the representatives
in (D.43) cancels the intermediate $u_\tau$, so these
are functors on all arrows and preserve every composition.

We specify why the retained operations also exist on this
index system.  Relabeling is the corresponding automorphism
of the surface group.  Filling one puncture is the quotient
of $\Pi_{g,n+1}$ by the closed normal subgroup generated by
its cusp.  An inner action before filling induces an inner
action afterwards, so the full forgetful homomorphism sends
$K_{g,n+1}$ into $K_{g,n}$.  Hence it is continuous for
the congruence topologies.  The transporter map is that
same homomorphism, with the source and target multicurves
stabilized; C.1 supplies its actual twist-lattice map.

For nodal clutching use the audited cutting theorem
\cite[Theorem 4.5, pp.~17--19]{ET_Boggi2014}.  Its pinned
stabilizer quotient by the full node-twist lattice is
exactly the product of the component procongruence
groups.  The quotient on the unpinned group is the actual
finite image $A_\sigma$ in the signed permutations of
the cut curves, preserving the labels and topological
type.  No split extension or larger permutation group
is substituted.  Thus the actual component clutching
maps are continuous for the congruence topologies.
On a retained-node diagram, cutting representatives
and extending them across its pinned collars are the
same inverse maps as in A.9--A.10; the displayed product
topology makes their transporter maps continuous too.
This covers self-clutching as well.  All comparisons of
E.39 therefore descend through the actual forgetful,
clutching and relabeling functors, and their coherence
equations remain equations after quotient and inverse
limit.  Genus zero has full congruence topology by D.19,
so restriction returns the given parameter.
\end{proof}

\begin{result}{D.32}{Proved}{The congruence kernel is the small-level fundamental group}
Let $d=3g-3+n$ and let $K_{g,n}$ be (D.40).  For every
stable type and every integer
\[
                       2\leq k\leq6,\qquad k<d,
\]
there is a natural isomorphism, after a basepoint choice,
\[
        \Pi_1(\check P_k)\ \cong\ K_{g,n}.                  \tag{D.44}
\]
Here $\check P_k$ is the inverse limit over congruence
levels of the rank-selected complex defined in D.5,
and $\Pi_1$ is Quick's profinite fundamental group.
The full level complex $P_k^{\mathrm{lev}}$ is its
connected simply connected profinite Galois covering,
with deck group $K_{g,n}$.

Consequently the assertion of procongruence simple
connectivity at any one of these levels is exactly
the vanishing of the corresponding congruence kernel.
At $k=2$ the completed local-group presentation of
D.11 still has total group $\widehat\Gamma_{g,n}$;
its map to $\check\Gamma_{g,n}$ has kernel (D.44).
For $g\leq2$ this kernel vanishes by D.19.  The first
ambient type with an undecided kernel here is $(3,0)$,
already at $k=2$.
\end{result}
\begin{proof}
Every vertex $\sigma$ of $P_k$ is nonempty because
$k<d$.  Cutting the connected surface along it gives
stable components, each with at least one node mark.
Their modular dimensions satisfy
\[
       \sum_v(3g_v-3+n_v)=d-|\sigma|\leq k.                  \tag{D.45}
\]
Each summand is nonnegative.  A component of genus
at least three has $n_v\geq1$ and dimension at least
seven.  Thus all cut components have genus at most
two in the stated range.

Write $C_\sigma$ for the actual discrete stabilizer,
and $C_{\vec\sigma}$ for its subgroup fixing the
curves with their orientations and all branch pins.
The full node-twist lattice has its full topology
in $C_{\vec\sigma}$ by E.29.  The discrete cutting
sequence and right exactness of completion therefore
give an exact top row in the diagram
\[
\begin{tikzcd}[column sep=small]
1\arrow[r]&\widehat{\mathbb Z}^{\sigma}\arrow[r]\arrow[d,equal]
 &\widehat C_{\vec\sigma}\arrow[r]\arrow[d]
 &\prod_v\widehat\Gamma_{g_v,n_v}\arrow[r]\arrow[d]
 &1\\
1\arrow[r]&\widehat{\mathbb Z}^{\sigma}\arrow[r]
 &\check C_{\vec\sigma}\arrow[r]
 &\prod_v\check\Gamma_{g_v,n_v}\arrow[r]&1.
\end{tikzcd}                                                   \tag{D.46}
\]
The bottom row and the actual closure
$\check C_{\vec\sigma}\subset\check\Gamma_{g,n}$
are precisely the valid Theorem 4.5 cited above.
D.19 makes the right vertical map an isomorphism.
The two exact rows then make the middle map an
isomorphism.  Restoring the same actual finite
permutation image $A_\sigma$ gives
$\widehat C_\sigma\cong\check C_\sigma$.
In particular $\widehat C_\sigma$ embeds into
$\widehat\Gamma_{g,n}$, because its composite into
$\check\Gamma_{g,n}$ is injective, and
\[
               K_{g,n}\cap\widehat C_\sigma=1.               \tag{D.47}
\]
The same holds for every flag stabilizer, which has
finite index in the stabilizer of its finest
multicurve.  This proves the exact local topology
needed here, rather than assuming it for arbitrary
full-completion stabilizers.

There are finitely many mapping-class orbits of
vertices and simplices.  On each orbit, the full
level completion is the homogeneous profinite set
$\widehat\Gamma/\widehat C$ and the congruence
completion is $\check\Gamma/\check C$, for its
actual flag stabilizer $C$.  Thus, degree by degree
and compatibly with every face and degeneracy map,
\[
 Y=P_k^{\mathrm{lev}},\qquad
                    \check P_k=K_{g,n}\backslash Y.
\]
Equation (D.47) proves that the continuous action
of $K_{g,n}$ on every vertex of $Y$ is free,
including profinite vertices: every such vertex
is a translate in one of these homogeneous orbits.

The quotient $Y\to\check P_k$ is an actual profinite
covering in the sense of
\cite[\S2.1, Propositions 2.2 and 2.27]{D_Quick2008}.
Indeed take any simplex downstairs and any prescribed
lift of one of its vertices.  Choose one lift of
the simplex.  There is a unique element of $K_{g,n}$
carrying its chosen vertex to the prescribed vertex,
by transitivity on that fibre and freeness.  Its
translate of the simplex is the required unique
lift.  This is exactly Quick's
$\Delta[0]\to\Delta[m]$ lifting definition of a
profinite covering.  The maps are continuous:
they are continuous bijections between compact
profinite spaces and Hausdorff spaces in each
degree.  The fibre is $K_{g,n}$ with its regular
action, so the covering is Galois.

D.11 proves that $Y$ is connected and has trivial
$\Pi_1$.  Quick's covering classification and the
exact sequence in the proof of Proposition 2.27
therefore give
\[
 1\longrightarrow\Pi_1(Y)\longrightarrow
 \Pi_1(\check P_k)\longrightarrow K_{g,n}\longrightarrow1.
\]
This proves (D.44) directly.  In particular no
unproved assertion that arbitrary homotopy groups
commute with inverse limits is used.  A choice of
lift of the base vertex identifies its fibre with
$K_{g,n}$; changing that choice conjugates the
identification, which is the stated basepoint
ambiguity.  Finally at $k=2$ the vertex and flag
groups in the full and congruence diagrams have
just been identified.  The completed presentation
of D.11 is therefore the same local-group
presentation, with its actual quotient map
$\widehat\Gamma\to\check\Gamma$.  Its kernel is
(D.44), as asserted.
\end{proof}
